\documentclass[aps,prl,reprint,amsmath,amssymb,nofootinbib,raggedbottom]{revtex4-2}
\usepackage[T1]{fontenc}
\usepackage{tnnotes}
\usepackage[colorlinks=true,citecolor=blue,linkcolor=blue,urlcolor=blue]{hyperref}
\hypersetup{pdftitle={How the Initial Estimate Is Certified}}

% ------------------------------------------------ diagrams specific to this note
% A periodic matrix-product vector against its copy: two rings of A tensors
% joined by physical legs; #1 = 1 inserts one bond operator h on two sites.
\newcommand{\mpsnorm}[1]{\tikz[tn,tn centered]{
 \foreach \y/\s in {.85/1,-.85/-1}{
   \draw[ring] (-.5,\y)--(-.5,{\y+\s*.75})--(6.7,{\y+\s*.75})--(6.7,\y)--cycle;}
 \foreach \x in {.5,2.0,4.6,6.0}{\draw[phys] (\x,-.85)--(\x,.85);}
 \tndots{3.3}{.85} \tndots{3.3}{-.85}
 \ifnum#1=1 \node[sym,minimum width=2.6em,minimum height=.7em,fill=mpoblue!20] at (1.25,0) {$h$};\fi
 \foreach \x in {.5,2.0,4.6,6.0}{
   \node[mpo,fill=white,minimum width=.9em] at (\x,.85) {$A$};
   \node[mpo,fill=white,minimum width=.9em] at (\x,-.85) {$A$};}}}
% Word-sized capped masses: bars at a dashed cap, the last one partial.
\newcommand{\capbars}{\tikz[baseline=0pt,x=.34em,y=1.1em]{
 \foreach \h [count=\i] in {.8,.8,.8,.35}{\fill[mpored!70] ({\i-1},0) rectangle ++(.78,\h);}
 \draw[black!55,densely dashed,line width=.3pt] (-.2,.8)--(3.95,.8);
 \draw[black!50,line width=.3pt] (-.2,0)--(3.95,0);}}
% Word-sized transfer spectrum at b=49/4: one bar above the cap l (red), the
% rest of J (red, light) and the three identical N lists (grey).
\newcommand{\seedspectrum}{\tikz[baseline=0pt,x=.34em,y=1.1em]{
 \fill[mpored] (0,0) rectangle ++(.78,1);
 \foreach \h [count=\i] in {.35,.2}{\fill[mpored!40] (\i,0) rectangle ++(.78,\h);}
 \foreach \h [count=\i from 3] in {.58,.58,.58,.3,.3,.3}{\fill[black!30] (\i,0) rectangle ++(.78,\h);}
 \draw[black!50,line width=.3pt] (-.2,0)--(9,0);}}

\begin{document}
\title{How the Initial Estimate Is Certified}
\date{October 15, 2026}
\begin{abstract}
The finite-scale input $\Ps(2304,784), \Pt(4608,784) \ge 0.99$ of Ref.~\cite{main} is established without direct tensor contraction at that scale.
Ref.~\cite{source} certifies a pair of bounds at length $36$ and $\beta=49/4$, which are then propagated to $(2304,784)$ through six scale doublings with decoupled defects.
The underlying verification reduces to three closed tensor networks: twisted rings of at most twelve sites, symmetry-projected rings, and an integer matrix-product trial state of bond dimension $26$ on $72$ sites.
\end{abstract}
\maketitle
\setlength{\abovedisplayskip}{5pt plus 2pt minus 1pt}
\setlength{\belowdisplayskip}{5pt plus 2pt minus 1pt}

\appendix
\setcounter{secnumdepth}{2}
\setcounter{section}{8}

\section{Architecture of the Certificate}\label{app:c-plan}
The proof of Ref.~\cite{main} takes one numerical input: the spatial and
thermal purities at length $2304$ and inverse temperature $784$ are at
least $0.99$. Computing them directly would require partition functions of
rings of $4608$ spins, far beyond any computer. Ref.~\cite{source}
certifies them instead from rings of at most twelve sites, and this note
explains how.

The idea is to start at a much smaller scale and let the doubling
recurrence carry the bounds up. At length $36$ and inverse temperature
$49/4$, Ref.~\cite{source} certifies upper bounds $p_0$ and $q_0$ on the
spatial and thermal defects. Six doublings take $(36,49/4)$ to
$(36\cdot2^6,\tfrac{49}{4}\cdot2^6)=(2304,784)$, and along the way both
defects fall below $1/100$.

We write $b$ for an inverse temperature in this note, mainly the certificate
value $b=49/4$, and $X_b$ for the spatial transfer operator of
Ref.~\cite{main} at $b$. Its eigenvalues $\lambda_i(b)$ are real, and
$Z_n(b)=\sum_i\lambda_i^n$. Following Ref.~\cite{source}, the Hamiltonian
is shifted by $\kappa=700741/500000$ per site, $H_n\to H_n+\kappa n$. This
$\kappa$ lies just below minus the ground-state energy per site (the trial
state of Sec.~\ref{app:c-spectrum} has energy per site below $-\kappa$), so
after the shift the leading eigenvalue of $X_b$ is close to one. The shift
cancels from every purity.

The certificate proceeds in four steps, from twisted partition functions
$\Tr(X_b^nU_g)$ to a leading-eigenvalue lower bound $\lambda_0>\ell$,
starting defects $(p_0,q_0)$ at $(36,49/4)$, and six doublings to the
starting estimate of Ref.~\cite{main}:
\begin{equation}
 \begin{aligned}
 \underbrace{\tnchain[node=xvertex,last=sym,ring]{,,\cdots,,U_g}}_{n\le12,\ b=49/4}
 &\ \xrightarrow{\ \text{sectors}\ }\
 \underbrace{\seedspectrum}_{\lambda_0>\ell}\\
 &\ \xrightarrow{\ \text{trial}\ }\ (p_0,q_0)
 \ \xrightarrow{\ \times6\ }\
 \underbrace{\Ps,\Pt\ge0.99}_{(2304,784)}.
 \end{aligned}
 \label{eq:c-plan}
\end{equation}
\begin{itemize}
\item \emph{Twisted partition functions} (Sec.~\ref{app:c-tori}): certified values of
$Z_n(b,g)$ for $4\le n\le12$, $b=49/4$, and two twists $g$.
\item \emph{Sectors and filters} (Sec.~\ref{app:c-spectrum}): bounds on
every eigenvalue of $X_b$, sector by sector. A trial vector then forces
one eigenvalue above $\ell=0.999$.
\item \emph{Initial defects} (Sec.~\ref{app:c-spectrum}):
$1-\Ps(36,b)\le p_0$ and $1-\Pt(72,b)\le q_0$.
\item \emph{Six doublings} (Sec.~\ref{app:c-doubling}): the coupled update
reaches $(2304,784)$.
\end{itemize}
Every number is an exact rational or an integer; floating-point
arithmetic is used nowhere in the proof.

\section{Twisted boundary conditions}\label{app:c-tori}
\emph{Twists and symmetries.} Untwisted traces $Z_n(b)=\sum_i\lambda_i(b)^n$
sum over all spin-rotation sectors at once; separating the $D_2$-invariant
sector from the three equivalent nontrivial sectors (the eigenvalue lists $J$
and $N$ of Sec.~\ref{app:c-spectrum}) requires twisted traces. A
proper signed permutation $g$ of the Cartesian axes acts on histories by $U_g$,
which relabels events and multiplies by their signs. Rotating every matrix
inside a site trace shows that $U_g$ slides through a transfer vertex,
\begin{equation}
 \tnchain[last=sym]{X_b,U_g}=\tnchain[node=sym,last=xop]{U_g,X_b},
 \label{eq:c-slide}
\end{equation}
so a single $U_g$ may be placed on any virtual bond of the ring. On the
physical side it conjugates one end of a single bond, the seam:
\begin{equation}
 Z_n(b,g)=\Tr X_b^nU_g
 =\tnchain[node=xvertex,last=sym,ring]{,,\cdots,,U_g}
 =\Tr e^{-bH_n(g)} .
 \label{eq:c-twist}
\end{equation}
The certificate needs $g=1$ and the $\pi$ rotation $g=P$.

\emph{Purification.} The physical side is computed as a squared
Frobenius norm, which splits the imaginary-time circle at half its length,
\begin{equation}
 Z_n(b,g)=\tncolumn[ring]{F,F^\dagger}=\|F\|_F^2,\qquad F=e^{-bH/2},
 \label{eq:c-cut}
\end{equation}
so that only columns $Fe_r$ are needed, one basis state $e_r$ at a time.

\emph{Uniformization.} Shift and scale $H$ so that
$R=1-(\tfrac32n+H)/16$ is a contraction. Since $8b=98$ at $b=49/4$,
uniformization (Jensen's method) writes $F$, up to the scalar
$e^{-b(\kappa-3/2)n/2}$, as a Poisson mixture of powers of $R$:
\begin{equation}
 F\propto e^{98(R-1)}=\sum_{j\ge0}e^{-98}\frac{98^j}{j!}\;
 \tncolumn[node=gibbs,gap=.15em]{R,R,\cdots,R}\,,
 \label{eq:c-poisson}
\end{equation}
with $j$ factors of $R$ in the $j$th term. This is the discrete counterpart of the
event expansion behind the thermal MPO: the events of $W_\beta$ arrive as a
Poisson process, and here each factor of $R$ is one event.

The sum is cut at $j=280$. The truncation error is below $2^{-84}$ in
operator norm, because $\|R\|\le1$~\cite{foxglynn}. The truncated sum is
evaluated by Horner's rule on integer vectors scaled by $Q=2^{55}$, taking
the floor after every multiplication by $R$. Each multiplication adds at most
$\sqrt{\dim}+1$ to the Euclidean error, and the contraction carries earlier
errors forward unamplified. Each column is therefore known to within
$281(\sqrt{3^n}+2)$ in units of $1/Q$.

\emph{Orbits.} The trace in Eq.~\eqref{eq:c-cut} needs one column per
basis word. Cyclic shifts, reversal and global spin flip permute the
columns without changing their norms. For the twisted ring, a shift moves
the seam, and a diagonal sign gauge moves it back. Only one representative
per orbit is computed, weighted by the orbit size. The results are eighteen
integers $W_{n,g}$ of about $30$ digits. Turning them into rational
enclosures gives the table of $Z_n(49/4,g)$, each value to within
$\pm3\times10^{-5}$. As an independent sanity check, not
part of the proof, floating-point exact diagonalization reproduces all
ten centers for $4\le n\le8$ to within $10^{-6}$.

\section{From partition functions to the transfer spectrum}\label{app:c-spectrum}
\emph{Sector projectors.} The four twists $g\in D_2=\{1,P_x,P_y,P_z\}$ commute
with $X_b$, so their character sums are projectors that can be inserted
into the ring:
\begin{equation}
 \begin{gathered}
 \Pi_\chi=\tfrac14\textstyle\sum_{g\in D_2}\chi(g)U_g,\\
 \Tr X_b^n\Pi_\chi=\tnchain[node=xvertex,last=sym,ring]{,,\cdots,,\Pi_\chi}.
 \end{gathered}
 \label{eq:c-proj}
\end{equation}
A cyclic rotation of the axes permutes the three nontrivial characters, so
those three sectors carry identical eigenvalue lists, called $N$. Write $J$
for the list of the invariant sector. Inverting the characters gives both
lists from the two twists alone:
\begin{equation}
 J_n=\tfrac14\bigl[Z_n+3Z_n(P)\bigr],\qquad
 N_n=\tfrac14\bigl[Z_n-Z_n(P)\bigr],
 \label{eq:c-sectors}
\end{equation}
with $Z_n=J_n+3N_n$. The coefficients sum to at most one in absolute value,
so each sector moment inherits the $3\times10^{-5}$ error. This holds for
odd $n$ too; odd moments carry sign information.

\emph{Polynomial filters} (the term of Ref.~\cite{source}). For any polynomial $F\ge0$ on $\mathbb R$, the closed
network $\Tr F(X_b)$ is a finite combination of ring partition functions:
\begin{equation}
 \sum_iF(\lambda_i)=\sum_jf_j\,\tnchain[node=xvertex,ring]{,\cdots,}_{\,j}
 \ \le\ B .
 \label{eq:c-filter}
\end{equation}
Every term on the left is nonnegative, so no single eigenvalue can have
$F(\lambda)>B$. Two degree-12 filters of the form
$x^4Q(x)^2$ give
\begin{equation}
 |\lambda|<v_J=1.00139\ \text{in }J,\qquad |\lambda|<v_N=0.872\ \text{in }N.
 \label{eq:c-caps}
\end{equation}
The filter needs $F\ge0$ on the whole real line, not only on
$[0,\infty)$, because $X_b$ is self-adjoint but not positive. Reflection
positivity is what makes the eigenvalues real~\cite{scope}.

\emph{Capped masses} (also Ref.~\cite{source}'s term; a convexity bound). The twelfth moments are bounded by
$m_J=1.0657205$ and $m_N=0.2783155$. Weights $t_i\le h$ with total mass at
most $m$ have their $r$th powers maximized by filling bars to the cap,
\capbars:
\begin{equation}
 \sum_it_i^r\le\mathcal C(m,h,r)=qh^r+(m-qh)^r,\qquad
 q=\lfloor m/h\rfloor .
 \label{eq:c-capped}
\end{equation}

\emph{The trial vector.} One matrix-product vector on $72$ sites, with real
integer tensors $A_s$ of bond dimension $26$, has energy per site below
$-\kappa$. Its norm and energy are two double-layer rings,
\begin{equation}
 \begin{gathered}
 \langle\psi|\psi\rangle=\mpsnorm0=\Tr T^{72},\\
 \langle\psi|H|\psi\rangle=72\,\mpsnorm1 .
 \end{gathered}
 \label{eq:c-trial}
\end{equation}
The virtual space carries four spin-$\frac12$, three spin-$\frac32$ and one
spin-$\frac52$ multiplets, $4\cdot2+3\cdot4+6=26$. The contraction splits
into blocks of dimension at most $162$, evaluated in exact integers. The
outcome, $E_0(72)<-72\kappa$, means that one Boltzmann factor of $Z_{72}(b)$
exceeds one, so
\begin{equation}
 Z_{72}(b)>1\quad\text{for every }b>0 .
 \label{eq:c-z72}
\end{equation}

\emph{A dominant eigenvalue.} Suppose every eigenvalue in $J$ had
$|\lambda|\le\ell=0.999$. The capped-mass bound applied to $J$ and to the
three copies of $N$ would give $Z_{72}(49/4)\le0.931<1$, contradicting
Eq.~\eqref{eq:c-z72}. Hence some $\lambda_0\in J$ has $|\lambda_0|>\ell$:
the spectrum looks like \seedspectrum, with one tall bar.

Removing $\lambda_0$ leaves the following bound on the $k$th moment of the rest:
\begin{equation}
 R_A(k)=3\,\mathcal C\bigl(m_N,v_N^{12},\tfrac k{12}\bigr)
 +\bigl(m_J-\ell^{12}\bigr)^{k/12}.
 \label{eq:c-rest}
\end{equation}

\emph{Initial defects.} Using $Z_{72}(b)\ge\lambda_0^{72}>\ell^{72}$ in
$\Ps(36,b)=Z_{72}(b)/Z_{36}(b)^2$ and $Z_{72}(2b)>1$ [Eq.~\eqref{eq:c-z72}] in
$\Pt(72,b)=Z_{72}(2b)/Z_{72}(b)^2$ bounds both purities at the base scale,
giving $(p_0,q_0)\simeq(0.0479,0.1815)$:
\begin{equation}
 \begin{aligned}
 \Ps(36,b)&=\frac{\Zt42}{\Ztsq21}\ge\Bigl(1+\frac{R_A(36)}{\ell^{36}}\Bigr)^{-2}
   =1-p_0,\\
 \Pt(72,b)&=\frac{Z_{72}(2b)}{Z_{72}(b)^2}\ge\bigl(v_J^{72}+R_A(72)\bigr)^{-2}
   =1-q_0 .
 \end{aligned}
 \label{eq:c-seed}
\end{equation}

\section{Six doublings}\label{app:c-doubling}
The thermal defect $q_0$ exceeds the diagonal threshold
$\varepsilon_*\simeq0.127$ of Ref.~\cite{scope}. A single common $\varepsilon$
would therefore not start the iteration. Instead, Moves~1--4 of Ref.~\cite{main}
track the two defects separately; with $f(z)=z^2/[2(1-z)^2]$ and
$\lceil\cdot\rceil_7$ rounding up to a multiple of $10^{-7}$,
\begin{equation}
 \begin{aligned}
 p_{j+1}&=\bigl\lceil f\bigl(1-(1-p_j)^2(1-q_j)\bigr)\bigr\rceil_7,\\
 q_{j+1}&=\bigl\lceil f\bigl(1-(1-q_j)^2(1-p_{j+1})\bigr)\bigr\rceil_7.
 \end{aligned}
 \label{eq:c-update}
\end{equation}

\begin{proofdiagram}
\begin{tikzpicture}[x=.95cm,y=1.55cm,font=\scriptsize]
 % log10 of the defects; data from Eq. (c-update), recomputed exactly
 \def\ps{-1.320,-1.218,-1.163,-1.147,-1.199,-1.425,-2.073}
 \def\qs{-.741,-.761,-.787,-.840,-.977,-1.351,-2.260}
 \draw[black!40,densely dashed,line width=.35pt] (0,-2)--(6,-2)
   node[right,black!60] {$10^{-2}$};
 \draw[black!40,densely dotted,line width=.35pt] (0,-.895)--(6,-.895)
   node[right,black!60] {$\varepsilon_*$};
 \foreach \v [count=\i from 0] in \ps{\ifnum\i>0 \draw[mpored,line width=.7pt] ({\i-1},\pl)--(\i,\v);\fi
   \fill[mpored] (\i,\v) circle[radius=1.3pt]; \xdef\pl{\v}}
 \foreach \v [count=\i from 0] in \qs{\ifnum\i>0 \draw[mpoblue,line width=.7pt] ({\i-1},\ql)--(\i,\v);\fi
   \fill[mpoblue] (\i,\v) circle[radius=1.3pt]; \xdef\ql{\v}}
 \node[mpored,anchor=north] at (2.5,-1.19) {$p_j$ (spatial)};
 \node[mpoblue,anchor=south] at (1.5,-.76) {$q_j$ (thermal)};
 \foreach \n [count=\i from 0] in {36,72,144,288,576,1152,2304}{
   \node[black!60] at (\i,-2.42) {$\n$};}
 \node[black!60,anchor=east] at (-.25,-2.42) {$n$};
\end{tikzpicture}
\end{proofdiagram}
\noindent The figure plots the defects on a logarithmic scale against the
length $n_j=36\cdot2^j$; the inverse temperature doubles alongside.
\begin{itemize}
\item The thermal defect falls from the first step.
\item The spatial defect first \emph{rises}, to $0.071$ at $j=3$: at small
$n$ the coupling identity costs more than squaring gains.
\item Both then collapse together and cross $10^{-2}$ at the last step,
\begin{equation}
 \begin{gathered}
 1-\Ps(2304,784)\le0.0084513,\\ 1-\Pt(4608,784)\le0.0054986 .
 \end{gathered}
 \label{eq:c-final}
\end{equation}
\end{itemize}
This recovers the initial estimate of Ref.~\cite{main} (Proposition~5.5 of Ref.~\cite{source}); re-evaluating the recurrence from the rational definitions of $p_0$ and $q_0$ confirms agreement digit for digit.

% Virtual-bond matrix gauge cancellation: W_j G_{j+1}^{-1} G_{j+1} W_{j+1}.
\newcommand{\gaugecancel}{\tikz[tn]{
 \draw (-.55,0)--(6.45,0);
 \foreach \x in {.35,5.55}{\draw[phys] (\x,-.8)--(\x,.8);}
 \node[mpo] at (.35,0) {$W_j$};
 \node[sym,minimum height=.85em,inner xsep=1.5pt] at (2.05,0) {$G^{-1}$};
 \node[sym,minimum height=.85em,inner xsep=2pt] at (3.75,0) {$G$};
 \node[mpo] at (5.55,0) {$W_{j+1}$};}}
\newcommand{\gaugedpair}{\tikz[tn]{
 \draw (-.55,0)--(3.15,0);
 \foreach \x in {.4,2.2}{\draw[phys] (\x,-.8)--(\x,.8);}
 \node[mpo] at (.4,0) {$\widetilde W_j$};
 \node[mpo] at (2.2,0) {$\widetilde W_{j+1}$};}}
\emph{The weaker initial estimate and temperature lifting.}
Ref.~\cite{source} also certifies $\Ps(60,105/4),\ \Pt(120,105/4)>7/8$ from a
table at $b=21/2$ for $n\le10$, three twists, the $N=120$ trial vector, and
$\sum_m z_m^{5/2}\le(\sum_m z_m)^{5/2}$ for the Boltzmann factors $z_m=e^{-bE_m}$,
giving $\Delta_L>\frac4{105}\log\frac{80}{79}$ for even $L\ge60$. Checking the
finite prefix intervals $[n_j,2n_j]$ via $\Pt(L,b_j)\ge(1-p_j)^2(1-q_j)$ and
sharpening $(v_J,v_N,\ell)$ in Eq.~\eqref{eq:c-seed} to
$(1.001351,0.871256,0.999998)$ already extends the $(36,49/4)$ certificate to
all even $L\ge36$ with $\Delta_L>0.01697$ (\path{models/spin1/experiments/}).
Raising the inverse temperature directly on the $b_0=49/4$ sector moments before
reducing to two purities gives a $19.1\times$ larger bound than Ref.~\cite{main}:
for even $12\le k\le60$ and $r\ge1$, Eq.~\eqref{eq:c-capped} bounds
$Z_k(b_0,P)\le J_k\le\mathcal C(m_J,v_J^{12},k/12)$ and
$Z_k(b_0)\le M_k:=\mathcal C(m_J,v_J^{12},k/12)+3\,\mathcal C(m_N,v_N^{12},k/12)$.
Raising the positive Boltzmann factors of the self-adjoint twisted Hamiltonians
to power $r$ gives $Z_k(rb_0)\le M_k^r=:D$ and
$Z_k(rb_0,P)\le\mathcal C(m_J,v_J^{12},k/12)^r=:E$, so at $b=rb_0$ the spatial
weights $u_i=|\lambda_i(b)|^k$ satisfy $\sum_i u_i\le D$, $J$-sector mass at
most $\tau=(D+3E)/4$, and $\sum_i u_i^{2h}=Z_{120}(b)>1$ ($h=60/k$). Because
each $N$ eigenvalue has multiplicity three, any $t>D/3$ with $Dt^{2h-1}<1$
forces $u_0\in J$ with $u_0>t$, leaving residual $s$th moment at most
\begin{equation}
 R_s(t)=\max\!\left\{3\Bigl(\frac{D-t}{3}\Bigr)^s,\
 (\tau-t)^s+3\Bigl(\frac{D-\tau}{3}\Bigr)^s\right\},
 \label{eq:c-resid-lift}
\end{equation}
upgraded whenever $(t')^{2h}+R_{2h}(t)<1$. For the physical Boltzmann factors
$z_m=e^{-b_0E_m(120)}$, $Z_{120}(b_0)\le M_{120}$ and $z_0>1$ give
$\sum_{m>0}(z_m/z_0)^r\le(M_{120}-1)^r$, yielding
\begin{equation}
 \begin{aligned}
 \Ps(60,rb_0)&\ge\Bigl(1+\frac{R_h(t)}{t^h}\Bigr)^{-2},\\
 \Pt(120,rb_0)&\ge\bigl(1+(M_{120}-1)^r\bigr)^{-2}.
 \end{aligned}
 \label{eq:c-templift}
\end{equation}
Combined with the $(36,49/4)$ prefix on $36\le L\le72$, Eq.~\eqref{eq:c-templift}
certifies $\Delta_L>0.07045>7/100$ at $(k,r)=(30,2)$ and $\Delta_L>0.07290$ at
$(28,17/8)$ for every even $L\ge36$. With a single untwisted bound
$Z_n(b_0)\le U$ and trial $Z_{2hn}(b_0)>1$ ($h,r\ge1$), the same argument with
$D=U^r$ and $t>D/2$ gives $\Ps(hn,rb_0)\ge[1+((D-t)/t)^h]^{-2}$ and
$\Pt(2hn,rb_0)\ge[1+(U^{2h}-1)^r]^{-2}$ whenever $t^{2h}+(D-t)^{2h}<1$
(\path{models/spin2/scripts/one_trace_target.jl}).

\section{Model-Dependent vs Universal Components}\label{app:c-general}
Each ingredient is a closed network, and only a few of them depend on the
model.
\begin{itemize}
\item \emph{Model independent:} the uniformization of Eq.~\eqref{eq:c-poisson},
the orbit reduction, the filters of Eq.~\eqref{eq:c-filter}, the capped
masses, and the doubling of Eq.~\eqref{eq:c-update}.
\item \emph{Model specific:}
 \begin{itemize}
 \item the symmetry group that splits the spectrum into sectors, which
 for spin~1 makes $N$ appear three times;
 \item one good trial vector;
 \item the sizes at which the filters isolate one eigenvalue.
 \end{itemize}
\end{itemize}

\emph{Smaller seed inputs for formal verification.}
Five studies in \path{models/lean_seeds} examine how to discharge the spin-1
seed inside the Lean~4 kernel (\path{lean/SeedKernel/}). First, even lengths
alone suffice: $n\in\{4,6,8,10,12\}$ with the two real twists $g\in\{1,P\}$ at
$b=49/4\to49/2$ and trace error $10^{-3}$ certifies $\Delta_L>0.039036$ for even
$L\ge60$, while $n\in\{6,8,10\}$ with twists $1,P,C$ ($C$ the $2\pi/3$ axis
permutation) at $b=21/2$ and error $3\times10^{-4}$ certifies
$\Delta_L>0.005966$ while cutting the orbit-trace dimension by $8.01\times$.
Second, replacing Poisson uniformization by low-energy trial eigenstates
\specbars{.92,.06,.02} (as done unconditionally for the two-state Ising seed in
Ref.~\cite{scope}) and a half-chain comparison operator
$H_n\succeq H_0=A_{n/2}\otimes I+I\otimes A_{n/2}+2cI$ (certified by exact
rational $81\times81$ block $\mathrm{LDL}^{\mathsf T}$ factorization) checks
$n=12$ packed \texttt{Nat} matrix-vector products in $9.6\text{--}26\text{ s}$.
Unlike the Ising seed, however, spin~1 at $n=12,b=49/4$ has $26$ ($g=1$) and
$34$ ($g=P$) eigenstates below the tail cutoff, and the $S^z=0$ comparison count
($49$ with a $32$-bin tail vs.\ $8$ physical levels) does not yet close without
homotopy continuation, so a kernel-checked spin-1 seed remains open.

\emph{Spin-2 trial certificates and the conditional one-trace target.}
For spin~1, shift-add bond kernels and decomposing the $D=26$ MPS transfer
matrix into six $\mathrm{SU}(2)$ multiplicity blocks (largest dimension $56$
vs.\ $162$ in $S^z$) speed up the exact length-$120$ trial contraction by
$360\times$ (\path{models/spin1}). For spin~2 ($\bm S^2=6I$, $5^{12}\simeq2.4\times10^8$
states, $\xi\approx49.5$), the reflection-positive bond form and the bootstrap
recurrence carry over unchanged, and $\mathrm{SU}(2)$-invariant MPS trials
certify $E_0(N)/N<-4.7595$ ($D=24$, integer) and $<-4.7612463$ ($D=295$, Arb;
reduced by Schur's lemma and the $\pi$-rotation involution from
$295^2=87{,}025$ to blocks of size at most $978$) up to $N=1024$
(\path{models/spin2}). By the untwisted one-trace bound above at
$(n,b_0,h,r)=(256,64,2,3)$ with shift $\kappa=4.7612463$, a single thermal upper
bound $Z_{256}(64)\le1.08$ would certify $\Delta_L>0.003019$ for even
$L\ge512$ (whereas the $(128,32)$ target at $\kappa=4.76122$ is ruled out by
the $N=128$ trial lower bound $Z_{128}(32)>1.1967$). Certifying
$Z_{256}(64)\le1.08$ remains open.

\emph{Certified thermal MPO compression.}
Two routes in \path{models/spin2} bound $Z_N(b)$ from above before MPO
compression (\path{models/spin2/IDEAS.md}): shifting $K=H_N+NI\succeq0$
($\|K\|\le5N$), a degree-$d$ Taylor majorant
$Q_d(K)=\sum_{j=0}^d(-\delta K)^j/j!+rI\succeq e^{-\delta K}$ ($\delta=b/(2m)$,
$m=2^p$) gives $Z_N(b)\le e^{(1-\kappa)bN}\Tr Q_d(K)^{2m}$, where the extensive
norm $\delta\|K\|\le5N\delta$ requires $p=28$ squarings at $N=256,b=64$ for
$d=2$, or $p=21$ for $d=4$ (factored into a positive rational sum of squares of
quadratics via the $\mathrm{LDL}^{\mathsf T}$ decomposition of the Laguerre
moment matrix, $P_4(x)=(1+\tfrac x2+\tfrac{x^2}{12})^2+\tfrac1{12}(x+\tfrac{x^2}2)^2+\tfrac1{72}(x^2)^2$);
alternatively, the Araki--Lieb--Thirring inequality for even and odd bond layers
$H_N=A+B$ gives $Z_N(b)\le\Tr[(e^{-bA/m}e^{-bB/m})^m]$ with step error governed
by the local bond norm rather than $N$. To certify MPO compressions along
repeated squarings, \path{CertifiedGaps.CertifiedMPO} (\path{julia/src/CertifiedMPO/})
uses two tensor-network mechanisms:
\begin{itemize}
\item \emph{Rational matrix gauges.} Diagonal virtual-bond rescalings leave the
local core-norm product $\prod_j\|W_j\|_{\mathrm{local}}$ at $N=256$ a factor of
$1.66\times10^{75}$ above the analytic upper bound $a_0$ on $\|Q_2/5^{N/2}\|$,
whereas invertible rational gauge matrices $\widetilde W_j=G_jW_jG_{j+1}^{-1}$
derived from the Schmidt decomposition cancel on every internal virtual bond,
\begin{equation}
 \gaugedpair\ =\ \gaugecancel\,,
 \label{eq:c-gauge}
\end{equation}
bringing the Cholesky-certified product
$\prod_j\|\widetilde W_j\|_{\mathrm{local}}$ within $1.000000701\,a_0$.
\item \emph{Three-layer trace correction and backward error propagation.} When a
squared step $A=B^*B$ is compressed to $C$ with $\|A-C\|_{\mathrm F}\le\eta$,
the identity
\begin{equation}
 \begin{aligned}
 \tncolumn[ring]{B^*,B,B^*,B}
 &=2\operatorname{Re}\tncolumn[ring]{C^*,B^*,B}-\tncolumn[ring]{C^*,C}\\
 &\quad+\|B^*B-C\|_{\mathrm F}^2
 \end{aligned}
 \label{eq:c-overlap}
\end{equation}
replaces the four-layer trace $\Tr(A^2)$ and linear triangle error
$(\|C\|_{\mathrm F}+\eta)^2$ by a three-layer signed overlap with quadratic
residual $\eta^2$ ($297\times$ tighter single-step excess at $D=32$). Across a
multi-step flow, propagating approximate backward adjoints (\path{AdjointTrace.jl})
retains the signed linear errors across steps ($1307\times$ narrower $N=16$
four-step enclosure), while backward Schatten bounds (\path{SchattenBridge.jl})
track $\|B_j\|_{p_j}$ with $p_j=2^{k+1-j}$ to use operator-norm residuals at
early steps and Hilbert--Schmidt residuals near the final trace. These
certificates pass replay on $4$- and $16$-site flows and the first two squarings
at $N=256$ ($b=2^{-20}$), leaving $26$ squarings to $b=64$ open.
\end{itemize}

\begin{thebibliography}{5}
\bibitem{main}
Companion note, \emph{A diagrammatic proof of the Haldane gap via spacetime
tensor networks}, \texttt{notes/thermal-mpo-gap-note.tex} (2026).
\bibitem{scope}
Companion note, \emph{Scope of the spacetime bootstrap},
\texttt{notes/thermal-mpo-gap-scope.tex} (2026).
\bibitem{source}
OpenAI, \emph{The periodic spin-one Haldane gap}, manuscript dated
September 24, 2026, Secs.~3--5 and Appendices A--C.
\bibitem{foxglynn}
B. L. Fox and P. W. Glynn, Computing Poisson probabilities,
Commun. ACM \textbf{31}, 440 (1988).
\end{thebibliography}
\end{document}
